	% ============================================================
	% 3. JUSTIFICACIÓN
	% ============================================================
	\section{Justificación}
	
La presente investigación se justifica por la necesidad de los establecimientos pequeños del sector retail peruano frente a las conductas asociadas al shoplifting, buscando desarrollar una alternativa eficiente para detectar rápidamente la presencia de hurtos en tiendas. En la actualidad, la mayoría de negocios dependen de soluciones de vigilancia convencionales, los cuales sufren de algunas limitaciones como: la fatiga visual, puntos ciegos y la dificultad para detectar patrones de comportamientos previos al hurto. A esto se suma que las soluciones basadas en aprendizaje profundo suelen requerir hardware costoso, lo que dificulta su adopción en este tipo de comercios.

Con respecto a la problemática mencionada, se propone una arquitectura de Deep Learning basado en la integración de la arquitectura YOLO con LSTM y el framework PCB (Pre Crime Behavior), enfocado a reconocer patrones espacio temporales en las conductas previas al crimen. La implementación estará enfocada a dispositivos de bajo consumo como NVIDIA Jetson o Raspberry Pi, buscando facilitar su aplicación en los comercios minoristas con bajo presupuesto. Esta propuesta se encuentra alineada con el objetivo general de desarrollar un modelo optimizado para mejorar la detección temprana de shoplifting en el contexto del retail minorista peruano.

Asimismo, la investigación busca aportar evidencia adaptada al contexto peruano por medio de la construcción de un dataset propio que incluya diversas secuencias de comportamientos sospechosos en diferentes condiciones de comercios minoristas locales. Esto permitirá considerar factores importantes como las variaciones de iluminación, ángulos de cámara, oclusiones y puntos ciegos.

Los resultados esperados ayudarán a determinar si la arquitectura propuesta puede lograr alcanzar un equilibrio entre la capacidad detección y la eficiencia computacional, proporcionando información sobre su viabilidad en la implementación para entornos de retail minorista peruano.
